\section{Conclusion}
\label{sec:conclusion}

Reactive temperature is an incomplete signal for tiering latency-sensitive
services because it identifies a working set only after that set becomes hot.
\system uses recurring execution phases to move a bounded set of pages before
their first access, while migration credits, leases, and an unchanged reactive
path contain prediction failures.  In our Linux prototype, this cross-layer
use of execution and memory signals reduces \ptail request latency by 31--58\%
without sacrificing throughput or isolation.  More broadly, CXL memory should
be managed not only by \emph{which} data is hot, but also by \emph{when} the
system can know that data will become hot.
